% LabLens — Complete Feature List
% AI-Powered Research Alignment Engine
% Columbia x NYU Claude Builder Hackathon — April 12, 2026

\documentclass{article}
\usepackage[margin=1in]{geometry}
\usepackage{enumitem}
\usepackage{titlesec}
\usepackage{hyperref}

\title{\textbf{LabLens — Feature List}}
\author{Columbia $\times$ NYU Claude Builder Hackathon}
\date{April 12, 2026}

\begin{document}
\maketitle

\section{Product Overview}

LabLens is an AI-powered research alignment engine that matches university students
to professors whose active, funded research aligns with their intellectual interests.
It goes beyond search or directory lookup — it is an intelligence layer that produces
actionable research guidance, personalized outreach, and privacy-preserving matching.

\section{Core ML Pipeline}

\begin{enumerate}[label=\textbf{\arabic*.}]

  \item \textbf{Intent Extraction (Claude)}
    \begin{itemize}
      \item Parses a student's free-text research interest into structured fields:
            primary domain, specific topics, methodological interests, career direction,
            search keywords, and an embedding-optimized description.
      \item Adapts extraction based on student level (undergrad, master's, PhD applicant).
    \end{itemize}

  \item \textbf{Live Faculty Discovery (Linkup API)}
    \begin{itemize}
      \item Multi-query agentic search to find professors at a given university
            in the student's research domain using live web data.
      \item Deep profile retrieval per professor: lab website, 5 most recent papers
            (2023--2025) with abstracts, active NSF/NIH grants, current PhD students
            and postdocs, recent graduates and their next positions, and open
            recruiting announcements.
      \item Lab alumni discovery via LinkedIn intelligence (PhD graduates,
            current positions, research topics).
    \end{itemize}

  \item \textbf{Semantic Embedding and Cosine Similarity Scoring}
    \begin{itemize}
      \item Uses OpenAI \texttt{text-embedding-3-small} (1536-dimensional vectors)
            to embed both the student's interest and each professor's paper abstracts
            and research statement.
      \item Computes cosine similarity between the student vector and each paper vector.
      \item Weighted scoring: 70\% from paper similarity (blend of max and average
            across up to 5 recent papers), 30\% from research statement similarity.
      \item Normalizes raw cosine (typically 0.1--0.9) to a 0--100 alignment score.
      \item Reports per-paper scores, best-matching paper index, and confidence level
            (high / medium / low based on number of papers embedded).
      \item Ranks all candidate professors by alignment score descending.
    \end{itemize}

  \item \textbf{Claude AI Synthesis Layer}
    \begin{itemize}
      \item Generates all qualitative intelligence via Anthropic Claude Sonnet,
            including professor cards, seed ideas, emails, skills gap analysis,
            lab culture reports, and timing assessments.
    \end{itemize}

\end{enumerate}

\section{Feature Modules}

\subsection{Professor Intelligence Card}
\begin{itemize}
  \item Plain-language summary of the professor's \emph{current} research focus
        (not career history).
  \item Three specific open research questions the professor is actively investigating.
  \item One-sentence explanation of why the student's interest aligns.
  \item Best paper to read first, with a reason why.
\end{itemize}

\subsection{Semantic Alignment Score and ML Explainability Panel}
\begin{itemize}
  \item Animated alignment score badge per professor (Strong Match $\geq$ 85,
        Good Match 70--84, Partial Match $<$ 70).
  \item Animated alignment bar on each professor card.
  \item Expandable ``How we scored this'' panel showing:
        model used, method (cosine similarity over research embeddings),
        raw cosine value, normalized score, and per-paper score breakdown
        with visual bars.
  \item Explicitly communicates that the score is a real ML computation,
        not an LLM opinion.
\end{itemize}

\subsection{Personalized Seed Research Ideas}
\begin{itemize}
  \item Three original research seed ideas per professor, each sitting at the
        intersection of the student's interest and the professor's open questions.
  \item Each idea includes: a title, a specific research question, connection to the
        professor's work, connection to the student's interest, and a difficulty
        rating (undergrad / master's / PhD level).
  \item Ideas are grounded in the professor's actual papers — never generic.
\end{itemize}

\subsection{Lab Culture Health Report}
\begin{itemize}
  \item Culture score (0--100) derived from publicly available signals.
  \item Evidence-based strengths and neutral-language watch-fors.
  \item Typical graduation timeline based on alumni data.
  \item Mentorship style signals inferred from public data.
  \item ``Best fit for'' student profile description.
  \item Three suggested questions to ask in an interview.
  \item Data sources: alumni time-to-graduation, student first-authorship rate,
        alumni destination diversity, lab size, public engagement, and public mentions.
\end{itemize}

\subsection{Lab Alumni Trail}
\begin{itemize}
  \item Tracks where past PhD graduates and postdocs went after leaving the lab.
  \item Shows each alumnus's name, graduation year, current role, employer,
        and research topics.
  \item Summary statistics: industry vs.\ academia placement count, average
        time in lab.
\end{itemize}

\subsection{Timing Intelligence Engine}
\begin{itemize}
  \item Determines the optimal timing for a student to contact a professor.
  \item Inputs: current date, semester and week, proximity to grant deadlines,
        grant deadline type, days since last arXiv submission, upcoming conference
        deadlines, and whether it is summer.
  \item Outputs: timing score (0--100), verdict (excellent / good / caution / wait)
        with color coding, primary reason, detailed explanation, specific optimal
        send time (e.g., ``Tuesday April 14 at 7:35 AM''), next contact window
        if verdict is ``wait'', and timing red flags.
  \item Calendar visualization showing semester position, grant deadlines,
        conference deadlines, and highlighted optimal contact windows.
\end{itemize}

\subsection{Skills Gap Analysis}
\begin{itemize}
  \item Compares the student's background against the professor's lab requirements
        (methods, programming tools, mathematical frameworks from papers).
  \item Lists required skills with importance level (required / helpful / bonus)
        and whether the student already has each.
  \item Identifies critical gaps and existing strengths.
  \item Provides a two-week prep plan with specific daily learning tasks and
        free resource links.
  \item Suggests how to frame the student's background in the email to
        minimize gap concerns.
\end{itemize}

\subsection{Grant Opportunity Matching}
\begin{itemize}
  \item Matches relevant fellowships (NSF GRFP, Hertz, Ford, etc.) to the student.
  \item Each match includes: name, amount, deadline, fit score, and
        ``Why this fits you'' text.
  \item Shows the professor's active grants that fund RAs (with amounts
        from NSF Awards API).
  \item Highlights when applying to a matched grant would strengthen an
        RA candidacy.
\end{itemize}

\subsection{Cold Email Generator}
\begin{itemize}
  \item Generates a personalized cold email draft per professor.
  \item Subject line references a specific paper title — never generic.
  \item Opening shows the student read the paper (references a specific finding
        or method).
  \item Body connects the student's background to a seed idea as an
        intellectual hook.
  \item Ask offers both a 15-minute call and email reply to reduce friction.
  \item Hard-capped at 180 words, enforced programmatically.
  \item Avoids clich\'{e} phrases (``I am reaching out'', ``Hope this finds you well'').
  \item Never starts with ``I'' — starts with something about the professor's work.
  \item Tone: peer-to-peer curiosity, not job application.
  \item Copy-to-clipboard button and ``Edit \& Personalize'' toggle.
\end{itemize}

\subsection{Follow-Up Tracker}
\begin{itemize}
  \item State machine: SENT $\to$ WAITING $\to$ FOLLOW\_UP\_READY $\to$
        FOLLOWED\_UP $\to$ REPLIED / ARCHIVED.
  \item Tracks email sent date, shows countdown to follow-up window (10 days).
  \item Auto-generates a follow-up email draft when the window opens.
\end{itemize}

\subsection{PhD Program Matching}
\begin{itemize}
  \item Separate flow triggered by ``I'm applying to PhD programs'' toggle.
  \item Ranks top 8 PhD programs using the same embedding pipeline.
  \item Per program: research fit score, number of active professors in the
        student's area, placement rate, average time-to-graduation, top 2
        professors to contact pre-application, and an application strategy note.
  \item Uses CSRankings data (via Linkup), NSF placement data, and semantic
        embeddings.
\end{itemize}

\section{Novel ML Contribution: Three-Tier Privacy-Preserving Matching}

\subsection{Problem Statement}
Professors often have confidential research projects (industry NDA, pre-publication,
grant-pending, politically sensitive) that they cannot publicly describe but for which
they need specific types of collaborators. No existing system enables matching across
this confidentiality boundary.

\subsection{Three-Tier Data Architecture}
\begin{enumerate}[label=\textbf{Tier \arabic*:}]
  \item \textbf{Public (weight 50\%)} — Published papers, abstracts, lab website,
        NSF/NIH grants, arXiv submissions, public engagement. Full Linkup search
        and full embedding.
  \item \textbf{Semi-Public (weight 30\%)} — Workshop papers, grant abstracts,
        course syllabi, GitHub repos, collaborator acknowledgments. Requires
        Linkup deep search to surface.
  \item \textbf{Private / Confidential (weight 20\%)} — Industry-funded projects,
        pre-submission work, unawarded grant proposals, sensitive topics. Professor
        inputs only ``what type of mind I need'' (skills and intellectual approach),
        \emph{not} what the project is. A differentially private embedding is
        generated and the raw text is immediately and irreversibly discarded.
\end{enumerate}

\subsection{Differential Privacy Mechanism}
\begin{itemize}
  \item Gaussian noise injection calibrated to a privacy budget $\varepsilon$
        (default $\varepsilon = 1.0$; options: 0.5 high privacy, 1.0 balanced,
        2.0 lower privacy / higher accuracy).
  \item L2 sensitivity of unit-normalized embeddings bounded by 2.
  \item Noise scale = $2 / \varepsilon$; noised embedding re-normalized to the
        unit sphere to preserve angular relationships for cosine similarity.
  \item Raw text is never stored, never logged, never returned — only the noised
        vector persists.
  \item Privacy middleware blocks request body logging on the private signal endpoint.
\end{itemize}

\subsection{Three-Tier Composite Scoring}
\begin{itemize}
  \item When a private signal exists: composite = $0.50 \times T1 + 0.30 \times T2
        + 0.20 \times T3$.
  \item When no private signal: weights renormalize to $0.625 \times T1 +
        0.375 \times T2$.
  \item Agreement bonus: +3 points if Tier 3 and Tier 1 scores agree within 15 points;
        $-2$ penalty otherwise (reflects noise uncertainty).
  \item Students see a higher alignment score with a lock badge and the message
        ``This professor has indicated additional research interest beyond their
        public profile that aligns with you'' — without ever seeing what was written.
\end{itemize}

\subsection{Professor Portal}
\begin{itemize}
  \item Dedicated \texttt{/professor} route for professors to manage their profile.
  \item Three-tier profile visualizer showing data coverage across tiers.
  \item Private signal input form with privacy level selector and explicit
        ``text discarded'' confirmation flow.
  \item Professors can delete their private signal at any time.
\end{itemize}

\section{User Interface and Pages}

\subsection{Page 1: Landing / Search}
\begin{itemize}
  \item Full-screen dark hero with animated dot grid background.
  \item Large textarea for natural-language interest input with rotating
        placeholder examples.
  \item University selector dropdown (Columbia, NYU, MIT, Stanford, CMU,
        UC Berkeley, Harvard — demonstrating generalizability).
  \item Student level selector (Undergrad / Master's / PhD Applicant).
  \item Feature badges: Live research data, ML alignment scoring, Linkup Search,
        Claude AI.
\end{itemize}

\subsection{Page 2: Results}
\begin{itemize}
  \item Left sidebar for filters and search refinement.
  \item Professor cards in a 2-column grid.
  \item Sort options: Best Match / Most Approachable / Newest Hire.
  \item Filter pills by department.
  \item Each card shows: alignment score badge, timing badge, current focus
        summary, animated alignment bar, seed idea count, alumni count,
        timing verdict, and grant count.
\end{itemize}

\subsection{Page 3: Professor Deep Dive}
\begin{itemize}
  \item Six tabbed sections, each lazy-loaded:
        (1) Research Intelligence — full card, seed ideas, best paper, ML score breakdown;
        (2) Lab Intelligence — culture report and alumni trail side by side;
        (3) Timing \& Approach — timing score, calendar visualization, optimal send time;
        (4) Skills Gap — lab requirements vs.\ student skills, two-week prep plan;
        (5) Grant Opportunities — matched fellowships and professor's active grants;
        (6) Your Email — generated email, copy button, edit mode, follow-up tracker.
\end{itemize}

\subsection{Page 4: PhD Program Matching}
\begin{itemize}
  \item Top 8 programs ranked by embedding-based research fit score.
  \item Per-program details and professor contact strategy links.
\end{itemize}

\subsection{Page 5: Professor Portal}
\begin{itemize}
  \item Three-tier data visualizer and private signal management interface.
\end{itemize}

\section{Loading, Error Handling, and UX Polish}

\begin{itemize}
  \item Five-step loading sequence with progress descriptions (searching faculty
        data, reading publications, computing alignment, analyzing culture,
        generating roadmap).
  \item Skeleton loaders with shimmer animation (no spinners).
  \item Graceful degradation: ``Limited public data'' badge when Linkup returns
        nothing; keyword-match fallback when embedding fails; refresh button
        when NSF API times out.
  \item Score numbers animate counting up on load; professor cards stagger
        in with 150ms delay; alignment bars animate width on mount; cards
        elevate with subtle glow on hover.
  \item Pre-loaded demo scenarios for reliable live demonstrations.
  \item Mobile responsive design.
\end{itemize}

\section{Tech Stack Summary}

\begin{itemize}
  \item \textbf{Frontend:} Next.js 14+ with TypeScript, Tailwind CSS, Framer Motion
  \item \textbf{Backend:} Python FastAPI
  \item \textbf{ML Core:} OpenAI text-embedding-3-small for semantic embeddings;
        NumPy for cosine similarity and Gaussian differential privacy
  \item \textbf{Search:} Linkup API (deep agentic search for live professor data)
  \item \textbf{AI Synthesis:} Anthropic Claude Sonnet for reasoning, summarization,
        seed idea generation, and email drafting
  \item \textbf{Public Data APIs:} NSF Awards API, arXiv API, academic calendar logic
  \item \textbf{Database:} In-memory for hackathon (no auth)
\end{itemize}

\end{document}
